\documentclass[10pt,a4paper]{amsart}
\usepackage[margin=19mm]{geometry}
\usepackage{amsmath,amssymb,mathtools}
\usepackage[scr=rsfs,cal=euler]{mathalfa}
\usepackage{xfrac}
\usepackage[unicode,hidelinks,bookmarks=false]{hyperref}
\newcommand{\ie}{i.e.\ }
\newcommand{\cf}{cf.\ }
\newcommand{\resp}{respectively\ }
\newcommand{\Subsection}{\subsection}
\providecommand{\nobreakdash}{\nobreak\hbox{-}}
\setlength{\emergencystretch}{2em}
\multlinegap0pt
\usepackage{bm}
\usepackage{bbm}
\usepackage{dsfont}
\usepackage{xspace}
\usepackage{cancel}
\usepackage{mathtools}
\usepackage{amsthm}
\makeatletter
\@ifundefined{thm}{\newtheorem{thm}{Thm}}{}
\@ifundefined{lem}{\newtheorem{lem}[thm]{Lemma}}{}
\makeatother
\newcommand{\NN}{\mathbb{N}}
\title[A congruence family modulo powers of 5 for generalized cubic]{A congruence family modulo powers of 5 for generalized cubic partitions via the localization method}
\author{Dalen Dockery}
\date{Source: arXiv:2508.05833v1 (CC BY 4.0)}
\begin{document}
\maketitle

\noindent\emph{Source and licence.} A congruence family modulo powers of 5 for generalized cubic partitions via the localization method. Version arXiv:2508.05833v1 (CC BY 4.0), distributed under the Creative Commons Attribution 4.0 International licence, as recorded in the arXiv licence metadata for this exact version and shown on the abstract page https://arxiv.org/abs/2508.05833v1. Full rights notice: licenses/arxiv-2508.05833-CC-BY-4.0.txt.\par\smallskip

\noindent\textit{Editorial scope.} Counted unit: the Lem whose source label is \texttt{lem:L\_seq} in the article (source0.tex lines 205--208, source label \texttt{lem:L\_seq}), delivered together with the article's own complete original proof of it (source0.tex lines 209--242, one \texttt{proof} environment of 34 lines). No mathematics is written, repaired or re-derived: statement and proof are the article's own text line for line. Required context retained uncounted: lem:U\_PS (lines 147--150). Every definition, display and label of those lines is retained unchanged. Omitted from this package and not redistributed: 1--146, 151--204, 243--692. Nothing inside a retained excerpt is omitted or altered: every retained body is delivered exactly as the article's lines are written, so no omission receipt of any kind accompanies this package. Citations: the retained excerpts carry no citation command at all, so no bibliography entry of the article is transcribed and no reference is redistributed. Reference adaptation: none was needed and none was made; every reference inside the delivered unit resolves inside it, so no unresolved reference or citation is produced. Printed slips kept literal: no printed word, symbol, hypothesis or number of the counted statement or of its proof was altered. Layout and class: the article's own class is \texttt{amsart} with packages including amssymb, amsmath, amsthm, cancel, hyperref, color, enumerate, verbatim, mathrsfs, mathtools, gensymb, caption, subcaption, graphicx; the wrapper is the lane's pinned \texttt{amsart} editorial wrapper at 10pt on A4, which restates the theorem-like environments the retained text needs and adds the article's own macro definitions that the retained text needs (\texttt{\textbackslash NN}), with extra wrapper packages (bm, bbm, dsfont, xspace, cancel, mathtools). The delivered numbering is the unit's own. Assets: no retained span contains a figure environment, a graphics-inclusion command, a PDF-page inclusion, a table or a TikZ picture; no image file is copied into this package, the package declares no asset dependency and no image file is redistributed with it. Licence: version arXiv:2508.05833v1 of this article is distributed under the Creative Commons Attribution 4.0 International licence, as recorded in the arXiv licence metadata for this exact version and shown on the abstract page https://arxiv.org/abs/2508.05833v1. A full rights notice is delivered at licenses/arxiv-2508.05833-CC-BY-4.0.txt. Characters: every retained excerpt is pure ASCII, so the delivered engine, XeLaTeX with its default Latin Modern text font, reports no missing character.
\begin{lem}\label{lem:U_PS}
Let $F(q)$ and $G(q)$ be $q$-series and take $d \in \NN.$ Then
\[U_d\left(F(q^d)G(q)\right) = F(q) \cdot U_d\left(G(q)\right).\]
\end{lem}

\begin{lem}\label{lem:L_seq}
For all $\alpha \geq 1$,
\[L_{2\alpha} = qf_1f_2^2 \sum\limits_{n=0}^\infty a_3\left(5^{2\alpha} n + \gamma_{\alpha}\right)q^n.\]
\end{lem}

\begin{proof} 
Proceed by induction. Note that 
\begin{align*}
L_1 &= U^{(0)}(1) = U_5(\Phi) \\
&= U_5\left(q^5 \frac{f_{25}f_{50}^2}{f_1f_2^2}\right) = qf_5f_{10}^2 U_5\left(\sum\limits_{n=0}^\infty a_3(n)q^n\right) \\
&= qf_5f_{10}^2 \sum\limits_{n=0}^\infty a_3(5n)q^n,
\end{align*}
thanks to Lemma \ref{lem:U_PS}. Thus
\begin{align*}
L_2 &= U^{(1)}(L_1) = U_5(L_1) \\
&=U_5\left(f_5 f_{10}^2 \sum\limits_{n=0}^\infty a_3(5n) q^{n+1}\right) \\
&= q f_1 f_2^2 \sum\limits_{n=0}^\infty a_3(25n+20) q^n.
\end{align*}
Now suppose that the conclusion holds for some $\alpha \geq 1.$ Then
\begin{align*}
L_{2\alpha+1} &= U^{(0)}(L_{2\alpha}) = U_5\left(q^6f_{25}f_{50}^2 \sum\limits_{n=0}^\infty a_3\left(5^{2\alpha} n + \gamma_{\alpha}\right)q^n\right) \\
&= qf_5f_{10}^2 U_5\left(\sum\limits_{n=0}^\infty a_3\left(5^{2\alpha}n+\gamma_\alpha\right)q^{n+1}\right) \\
&= qf_5f_{10}^2 \sum\limits_{n=1}^\infty a_3\left(5^{2\alpha+1}n - 5^{2\alpha}+\gamma_\alpha\right) q^n, \\
\end{align*}
again thanks to Lemma \ref{lem:U_PS}. Hence
\begin{align*} 
L_{2(\alpha+1)} &= L^{(1)}(L_{2\alpha+1}) = U_5(L_{2\alpha+1}) \\
&=U_5\left(qf_5f_{10}^2 \sum\limits_{n=1}^\infty a_3\left(5^{2\alpha+1}n - 5^{2\alpha}+\gamma_\alpha\right) q^n\right) \\
&= f_1f_2^2 U_5\left(\sum\limits_{n=2}^\infty a_3\left(5^{2\alpha+1}n - 5^{2\alpha+1} - 5^{2\alpha}+\gamma_\alpha\right) q^n \right) \\
&= qf_1 f_2^2 \sum\limits_{n=0}^\infty a_3\left(5^{2(\alpha+1)}n + 5^{2\alpha+2} - 5^{2\alpha+1} - 5^{2\alpha} + \gamma_\alpha\right)q^n \\
&= qf_1f_2^2 \sum\limits_{n=0}^\infty a_3 \left(5^{2(\alpha+1)}n + \gamma_{\alpha+1}\right) q^n,
\end{align*}
as
\begin{align*}
5^{2\alpha+2} - 5^{2\alpha+1} - 5^{2\alpha} + \gamma_\alpha &=  19 \cdot 5^{2\alpha} + \left(20 + \frac{19 \cdot 25(25^{\alpha-1}-1)}{24}\right) \\
&= 20 + \frac{19 \cdot 25(25^{\alpha}-1)}{24} \\
&= \gamma_{\alpha+1}.
\end{align*}
\end{proof}

\end{document}
